\subsection{Foundational Scaling Laws of Neural Language Models}
The study of scaling laws has become central to understanding and advancing large language models (LLMs). Kaplan et al.~\citep{kaplan2020scaling} first demonstrated that language modeling loss follows a power-law with respect to model parameters ($N$), dataset size ($D$), and training compute ($C$), providing a quantitative foundation for predicting performance at scale. Hoffmann et al.~\citep{hoffmann2022training} build on this foundation by focusing on compute-optimal training under a fixed FLOP budget, demonstrating that both model size \(N\) and training tokens \(D\) should scale equally to maximize performance. Pearson and Song~\citep{pearce2024reconciling} found much of the discrepancy stems from Kaplan’s focus on non-embedding parameters at small scales, while Porian et al.~\citep{porian2024resolving} highlighted the roles of last-layer costs, optimizer warmup, and scale-specific tuning in aligning the two laws. Furthermore, efforts to refine the scaling paradigm include exploring performance gains via data pruning~\citep{sorscher2022beyond}, the effects of repeated data~\citep{hernandez2022scaling}, and the predictive value of gzip-based data complexity measures~\citep{pandey2024gzip}. Extensions also examined scaling under data constraints~\citep{muennighoff2023scaling}, synthetic data regimes~\citep{qin2025scaling}, and downstream translation tasks~\citep{isik2024scaling}. Snell et al.~\citep{snell2024scaling} further demonstrate that increasing test-time compute can yield gains comparable to enlarging model size, pointing to an inference-time analogue of conventional training-time scaling laws.

\subsection{Reinforcement Learning Scaling Dynamics in Large Language Models}
Scaling laws in reinforcement learning have revealed systematic power-law relationships between model size, compute, and performance in interactive settings \citep{hilton2023scaling}. In the context of large language models (LLMs), works including DeepSeek-R1 \citep{guo2025deepseekr1}, KIMI K1.5 \citep{kimi2025scaling}, and the Agent RL Scaling Law \citep{mai2025agent} demonstrate that scaling reinforcement learning fine-tuning (RFT) along compute and training horizons leads to systematic improvements in reasoning, especially for mathematics and code. Prolonged training regimes \citep{liu2025prorl} and even one-example reinforcement learning \citep{wang2025reinforcement} show that extended or extreme low-data RL can expand reasoning boundaries, while minimal-data paradigms like LIMR highlight efficiency tradeoffs \citep{li2025limr}. Yet the dynamics remain fragile: SimpleRL-Zoo identifies optimization instability in open-base models \citep{zeng2025simplerlzoo}, and critical analyses question whether RL consistently provides reasoning gains beyond the base model \citep{yue2025doesrlincentivize}. Beyond compute and optimization, data composition also plays a central role: multi-domain studies reveal both synergy and interference across math, code, and logical reasoning, underscoring the importance of curation and initialization \citep{li2025onedomainhelps, cheng2025revisiting}. Recent efforts further improve RL-based reasoning through attention-guided process supervision~\citep{nie2026attnpo}, self-search reinforcement learning~\citep{fan2025ssrl}, and structured analysis of LLM reasoning behaviors~\citep{yuan2025understandingllmreasoningabstractive}. Taken together, these works suggest that RL scaling has the potential to expand reasoning capabilities in LLMs, but its governing laws remain unsettled, motivating our focused study in mathematical reasoning.


\subsection{Scaling Behaviors in Mathematical Reasoning with LLMs}
Beyond general and RL-specific scaling laws, several studies have targeted mathematical reasoning as a domain where scaling effects are most pronounced. Scaling Relationship on Learning Mathematical Reasoning with LLMs \citep{touvron2023scaling} analyzes performance across model families, showing that accuracy improves with scale while emergent verification behaviors remain inconsistent. Skywork-Math \citep{ye2024skywork} establishes data scaling laws, demonstrating that both the quantity and quality of mathematical corpora critically shape scaling curves. MAMMOTH \citep{yue2023mammoth} develops math-generalist models through multi-task training, revealing divergent scaling dynamics between specialist and generalist setups. Most recently, the Agent RL Scaling Law \citep{zeng2025agent} uncovers novel dynamics introduced by reinforcement learning with code execution, where emergent behaviors such as spontaneous tool use arise predominantly in mathematical problem solving. Collectively, these works suggest that scaling in mathematical reasoning is shaped by model size, data distribution, and training paradigm, while the underlying principles remain not fully understood.
